% --- PACOTES ---
%> Pacotes essenciais
\usepackage{graphicx}
\usepackage{indentfirst}
\usepackage[utf8]{inputenc}
\usepackage[T1]{fontenc}
\usepackage[brazilian]{babel}

%> Pacotes de estilo
\usepackage{titlesec}
\usepackage{multicol}
\usepackage{epigraph}
\usepackage{capt-of}
\usepackage{microtype}
\usepackage[autostyle]{csquotes}
\usepackage{booktabs}
\usepackage{enumitem}
\usepackage{array}
\renewcommand{\arraystretch}{1.6}
\usepackage[svgnames]{xcolor}
\usepackage{pdfpages} %> Necessário para adicionar páginas em pdf (registro, etc...)
\usepackage{fancyhdr}
\setlength{\headheight}{14.5pt}
%\usepackage{markdown} %> Para escrever linhas de código

%> Pacotes matemáticos
\usepackage{amsmath, amssymb, amsfonts}
\usepackage{esint}
\usepackage{upgreek}
\usepackage{tikz}    
\usetikzlibrary{angles, quotes, arrows.meta, babel}
\usepackage{xsim}
\usepackage[most]{tcolorbox}

%> Configurações das páginas
\pagestyle{plain}

%> Gerador de blá blá blá
\newcommand{\blablabla}[1][1]{%
  \foreach \n in {1,...,#1}{%
    Neste parágrafo vamos simular um texto em português para testar a tipografia, o espaçamento, a hifenização e a formatação geral do documento. Como estamos desenvolvendo um modelo robusto para a Universidade Federal de Viçosa (UFV), é importante que as margens sejam testadas com palavras reais da nossa língua, contendo acentuações, cedilhas e construções frasais típicas do nosso idioma. Dessa forma, garantimos que o código do template funcione perfeitamente, sem depender de pacotes externos ou línguas mortas. \par\vspace{0.2cm}%
  }%
}